\section{Why the five rules recover full intuitionistic logic}
\label{sec:counter-logic}

The clause method, gate construction, and logical completeness argument in
this section are inherited from the ConciseIPRA developments \cite{basis,counter,single}. We give
the details for the present counter-aware terms, including the
witness-sensitive introduction. The immediate predecessor $P(Sn)=n$
also shortens the conjunction and universal macros: no parity iterator is
needed. All statements below concern arbitrary retained contexts.

\subsection{The clause interface and its first admissible operations}

Formulas are $F::=(s=t)\mid D_{x,y}F$. A context is a reusable list of
clause occurrences $[A;y.B(y)]$, where $y$ is absent from $A$. Put
$H=(0=0)$. An ordinary assumption $F$ means the constant clause
$[H;v.F]$, with $v$ unused. In addition to Refl and counter evaluation
CE, the logical instructions are exactly
\[
 \frac{\Gamma\vdash A}{\Gamma\vdash B(t)}\ \mathrm{Use}
 \quad([A;y.B]\in\Gamma),
\]
\[
 \frac{\Gamma\vdash D_{x,y}B(x,y)\quad
       \Gamma,[B(b,0);y.B(b,y)]\vdash F}
      {\Gamma\vdash F}\ \mathrm{Open},
\]
\[
 \frac{\Gamma,B(t,\eta)\vdash B(t,Sz)}
      {\Gamma\vdash D_{x,y}B(x,y)}\ \mathrm{Close},
 \qquad
 \eta=\begin{cases}
 z,& B\text{ is a literal equation and }t\equiv0,\\
 0,&\text{otherwise}.
 \end{cases}
\]
Here $b$ in Open is fresh for the retained context, conclusion, and passive
matrix parameters. In Close, $z$ is fresh for the context, witness, and
passive matrix parameters. All substitutions are capture avoiding and
canonical; $\equiv$ tests literal syntax. An equation under a $D$ is not
a literal equation. Local assumptions and opened clauses are discharged.

Refl proves $H$, so Use gives ordinary assumption use. Weakening,
injective fresh renaming, and ordinary-assumption cut follow directly by
induction on finite derivations. For cut, replace every Use selecting the
distinguished constant clause $[H;v.F]$ by the supplied proof of $F$;
the proof of its antecedent $H$ can be discarded. Freshen eigenparameters
before inserting the proof. These operations preserve the literal-zero
witness test. Arbitrary numerical substitution requires the repair below;
it is not assumed at this stage.

\subsection{Cases and induction, before compound formula transport}

The arithmetic equality lemma from Refl and CE gives transport in every
numerical template and hence every equation template. Put
\[
                 Z=R(S0,0;i,u.0).
\]
Its syntax is not the literal zero, while successor computation proves
$Z=0$. Suppose an equation matrix $E$ has a case branch
$\Gamma,E(0,0)\vdash E(0,Sz)$. Under the assumption $E(Z,0)$,
equation transport gives $E(0,0)$; cut inserts the supplied branch; reverse
transport gives $E(Z,Sz)$. Close at witness $Z$ is in case mode and
concludes $D_{x,y}E(x,y)$. This proves every missing equation-case instance
at witness zero. Every other case instance was already primitive. Thus
\[
 \frac{\Gamma,B(t,0)\vdash B(t,Sz)}
      {\Gamma\vdash D_{x,y}B(x,y)}\quad\text{(Case)}
\tag{Case}
\]
is admissible for every matrix and every witness. Its repair used only
atomic equality transport, not transport through a binder.

For an unused coordinate $u$, abbreviate $D_{u,y}A(y)$ temporarily by
$q_yA$. Case introduces it from $A(0)\vdash A(Sz)$. Open and Use
specialize it: from $q_yA$ and $A(0)$ derive $A(t)$.
Consequently the following two derived rules hold:
\[
 \frac{\Gamma\vdash F(0)\quad\Gamma\vdash F(Sz)}
      {\Gamma\vdash F(t)}\quad\text{(numerical cases)},
 \qquad
 \frac{\Gamma\vdash E(0)\quad\Gamma,E(z)\vdash E(Sz)}
      {\Gamma\vdash E(t)}\quad\text{(equation induction)}.
\tag{NC, EI}
\]
In the first rule, Case introduces $q_yF$ while ignoring its assumed base.
In the second, Close at unused witness $0$ introduces $q_yE$ directly
in progress mode. In both cases Open supplies the conditional clause,
the given base supplies its antecedent, and Use at $t$ finishes. The
displayed $z$ is fresh for the retained context and passive parameters.
Only $E$ is required to be a literal equation.

Capture-avoiding numerical substitution is now admissible by proof
induction. A progress introduction keeps its equation root and literal-zero
witness. A case introduction whose witness becomes zero is replaced by
the derived Case rule. Refl, CE, Use, and Open commute with canonical
substitution, using Lemma~\ref{lem:canonical} and
freshening eigenparameters first. This justifies the substitutions of
derivations used below without assuming witness-test invariance.

\subsection{Transport, gates, and selection}

Full formula transport is proved by external induction on $F(v)$.
Equations were already handled. For $F(v)=D_{x,y}B(v,x,y)$, open
$F(s)$ at fresh $b$, retaining $[B(s,b,0);y.B(s,b,y)]$.
Introduce $F(t)$ by Case at witness $b$. Convert the local assumption
$B(t,b,0)$ backwards along $s=t$ using the smaller-template induction
hypothesis; Use the old clause at $Sz$; convert its result forwards to
$B(t,b,Sz)$. Close Case and Open. This proves
\[
                  s=t,\ F(s)\vdash F(t).
\]
The same construction proves binder congruence for uniform pointwise
interderivability: convert the target base backwards, specialize the old
clause, and convert its successor instance forwards. It uses Case and cut,
not induction on a compound arithmetic formula.

The previously computed selector satisfies $C(0,a,b)=a$ and
$C(Sz,a,b)=b$. Define gates on finite formula syntax by
\[
 G_v(s=t)=(C(v,0,s)=C(v,0,t)),\qquad
 G_v(D_{x,y}B)=D_{x,y}G_vB,
\]
with binders renamed away from $v$. Atomic equality and binder congruence
prove, by external formula induction,
\[
                 \vdash G_0F,\qquad G_{Sz}F\dashv\vdash F.
\]
For the first statement at a paired formula, use Case and the already
proved gate theorem for its smaller matrix; the assumed base is unused.

Padding $L(F)=D_{u,y}G_yF$, with fresh unused $u,y$, is interderivable
with $F$. Forward, Case uses $F$ to prove the successor gate. Backward,
Open, the zero-gate theorem, and Use at $1$ give $G_1F$, hence $F$.
Every formula is a finite chain of paired binders ending in an equation.
Pad the shorter of $A,B$ to equalize their depths and align bound names.
At their atomic leaves set
\[
 M_v((s=t),(a=b))=(C(v,s,a)=C(v,t,b))
\]
and carry the matched binders outside. Thus
\[
                M_0(A,B)\dashv\vdash A,\qquad
                M_{Sz}(A,B)\dashv\vdash B.
\tag{Selector endpoints}
\]
All these are finite macros in the original Eq/$D$ language.

\subsection{Every ordinary intuitionistic connective and quantifier}

With all auxiliary binders fresh, and $u$ unused, define
\[
\begin{aligned}
 \bot_*&=(0=1),\\
 A\to_*B&=D_{u,y}M_y(A,B),\\
 A\land_*B&=D_{u,y}G_yM_{P(y)}(A,B),\\
 \exists_*x\,A(x)&=D_{x,y}G_yA(x),\\
 A\lor_*B&=\exists_*x\,M_x(A,B),\\
 \forall_*x\,A(x)&=D_{u,y}G_y A(P(y)).
\end{aligned}
\tag{Logical macros}
\]
Here $P(Sz)=z$, so $P(1)=0$ and $P(2)=1$. Its value $P(0)=1$
causes no problem: the outer gate makes the zero instance provable.

\paragraph{Implication.}
Given a derivation of $B$ under ordinary assumption $A$, Case introduces
$A\to_*B$: the selector's zero instance gives $A$, and its successor
instance is equivalent to $B$. For elimination, Open retains the selector
clause; the supplied $A$ gives its zero antecedent and Use at $1$ gives
$B$. Discharge and eigenvariable conditions are exactly the ordinary ones.

\paragraph{Conjunction.}
From $A$ and $B$, numerical cases prove $M_v(A,B)$ for arbitrary $v$.
Substitute $v=P(Sz)$, apply the successor gate, and use Case to introduce
$A\land_*B$. Its zero antecedent is a gate theorem. For elimination,
Open and Use at $1$ and $2$ give the selector at $P(1)=0$ and
$P(2)=1$, hence $A$ and $B$. No parity or range theorem is used.

\paragraph{Existential quantification.}
From $A(t)$, Case at witness $t$ introduces its gated paired formula.
Conversely Open at fresh $b$, prove the zero-gate antecedent, and Use at
$1$ to obtain $A(b)$. Cut this into any supplied branch whose conclusion
and retained context are free of $b$, then close Open.

\paragraph{Disjunction.}
Choose witness $0$ or $1$ in the existential macro to obtain the two
injections. For elimination, first discharge the two supplied branches
as $A\to_*F$ and $B\to_*F$. Numerical cases, in the original context,
prove $M_v(A,B)\to_*F$ for arbitrary $v$, using the selector endpoints.
Existential elimination supplies $M_b(A,B)$ at a fresh tag $b$;
implication elimination gives $F$. This construction does not substitute
the tag inside a retained assumption: numerical cases were performed
before that assumption was added.

\paragraph{Universal quantification.}
For introduction, a derivation of $A(z)$ with $z$ fresh for the context
gives the Case branch $G_{Sz}A(P(Sz))$, by $P(Sz)=z$ and the gate
equivalence. Introduce the unused witness to obtain $\forall_*xA(x)$.
For elimination, Open its clause and prove its zero-gate antecedent.
Use at $St$ gives $G_{St}A(P(St))$, hence $A(t)$. Thus every numerical
target is covered directly; no claim about $P(0)$ is needed.

\paragraph{Falsity and numerical coverage.}
Under $0=1$, transport through the term $C(v,s,t)$ gives every equation
$s=t$ from its two selector computations. External induction on formula
depth, using Case, gives every Eq/$D$ formula. This is ex falso.
Numerical cases on
$n=0\lor_*\exists_*m\,(n=Sm)$ prove coverage: use the left injection
at zero, and the right injection with witness $z$ at $Sz$.
Coverage can then be used in contexts already containing $n$, through
ordinary disjunction elimination and equality transport.

\subsection{Faithful ordinary realization and both formula composites}

Let $T_R$ be ordinary intuitionistic first-order logic with equality over
\emph{raw} $R$ terms, coverage, separation $0=1\to\bot$, contextual
literal-equation induction, and both computation schemas
\[
 R(0,a;i,u.h)=Sa,\qquad
 R(Sn,a;i,u.h)=h[n/i,R(n,a;i,u.h)/u].
\]
Canonical native terms are read as these raw terms. The semantic
justification of canonical substitution uses the ordinary theory's
body-extensionality argument in Section~\ref{sec:normalization-soundness};
it is separate from the elementary syntactic coherence lemma. Read
\[
 (D_{x,y}B)^\circ=\exists x(B^\circ(x,0)\to\forall yB^\circ(x,y)),
 \qquad [A;y.B]^\circ=A^\circ\to\forall yB^\circ(y).
\]
Refl and CE are sound by the arithmetic equality and canonical-substitution
lemmas. Use and Open are ordinary implication/universal and existential
elimination. Close in case mode is sound by coverage. In progress mode,
assume $E(0,0)$ and use equation induction with its premise to prove
$\forall yE(0,y)$, then choose witness $0$. Thus every primitive native
derivation has its ordinary realization in $T_R$.

Conversely the derived logical rules above translate every ordinary
intuitionistic inference. Numerical equality and both computations are
derived from Refl and CE; coverage and equation induction were derived
above. The separation axiom translates to the immediate implication
$(0=1)\to_*(0=1)$. Hence all $T_R$ derivations translate into the native
calculus. The ordinary reading of each displayed macro is provably
equivalent to its ordinary connective: use the gate and selector endpoints,
coverage for arbitrary controls, and $P(St)=t$ for the universal macro.
External induction on ordinary formulas proves this formula round trip.

For the other composite, externally induct on a native formula. Equations
are unchanged. For $D_{x,y}B$, first use the induction hypothesis for its
matrix uniformly. From the original paired formula, Open at $b$. Under
its base assumption, Use its clause at a fresh $z$ to obtain $B(b,z)$;
universal introduction, implication introduction, and existential
introduction then build the macro translation of
$\exists x(B(x,0)\to\forall yB(x,y))$. Close Open.
Conversely, existential elimination on that macro gives a fresh $b$ and
the implication from $B(b,0)$ to its universal closure. Introduce the
original paired formula by Case at $b$: its local base feeds the implication,
and universal elimination at $Sz$ supplies the required successor instance.
The matrix induction equivalences convert the occurrences in both directions.
This proves native interderivability with its round trip under arbitrary
ordinary side assumptions. Both numerical maps in this logical comparison
are identities; no extra domain or equality quotient is introduced.

The separation clause in the ordinary realization is explicit. Merely
collecting the positive Eq/$D$ theorem images would admit singleton models
and would not be the asserted ordinary theory. The argument establishes the
full first-order consequence relation; it neither postulates excluded middle
for quantified formulas nor introduces induction for compound native formulas.
